\chapter{What the Next Decade Will Decide}
\label{chap:nextobs}

The signatures of Chapter~\ref{chap:nowobs} are already in hand, if ambiguous.
This chapter looks one step out, to predictions sharp enough that planned
programs---the next generation of CMB experiments, gravitational-wave satellites,
dark-matter detectors, and galaxy surveys---should decide them within five to
fifteen years. Several are clean discriminators against $\Lambda$CDM and the
Standard Model; one or two could, on their own, strongly favour the framework or
strongly embarrass it.

\section{The microwave background, in finer detail}

If the background is the parent's Hawking glow (Chapter~\ref{chap:home}), its
spectrum should be a perfect Planck curve---which it is, to one part in
$10^{5}$---and its departures from perfection should carry the parent's
fingerprints. Three are worth looking for: faint spectral features from the
membrane geometry; a correlation between the background temperature and the
dark-matter distribution (both being projections of the same parent through the
same surface); and specific large-angle polarization patterns from a rotating
parent, aligned with the preferred axis of Chapter~\ref{chap:nowobs}. CMB-S4 and
LiteBIRD will tighten the spectrum and polarization by orders of magnitude. The
harder, prior task is internal: re-deriving the acoustic-peak structure from the
bounce rather than from recombination, which standard cosmology fits superbly and
which the framework still owes---the central calculational challenge of
Chapter~\ref{chap:sims}.

\section{The neutrino sector's unmatched number}
\label{sec:lepton-asymmetry}

The matter excess $\eta\simeq6\times10^{-10}$ is among the best-measured numbers
in cosmology (Chapter~\ref{chap:nowobs}); its lepton-sector counterpart---the net
neutrino--antineutrino excess of the relic neutrino background---is almost
entirely \emph{un}measured, and that gap is a sharp test. Standard cosmology
usually \emph{assumes} the lepton asymmetry is comparable to the baryon one, but
the data only bound the common neutrino degeneracy to
$|\xi_\nu|\lesssim0.04$~\citep{oldengott2017}---a ceiling that would permit a relic
asymmetry up to $\sim\!10^{7}$ times the baryon asymmetry, a vast unprobed window.

The framework predicts that window is \emph{empty}. The matter bias is inherited
and then processed at the bounce by electroweak sphalerons, which are active there
($T\sim10^{7}\,\mathrm{GeV}$, far above their $\sim\!130\,\mathrm{GeV}$ freeze-out;
Chapter~\ref{chap:dawn}). Sphalerons in equilibrium conserve $B-L$ but lock the
baryon and lepton numbers into a fixed ratio---for the Standard Model
$|L/B|=\tfrac{51}{28}\simeq1.8$---so the relic neutrino degeneracy is
$\xi_\nu\sim\eta\sim4\times10^{-9}$, about $10^{7}$ below the current ceiling and
indistinguishable from zero. There is no large primordial lepton asymmetry to
spend.

\paragraph{The test.} A \emph{measured} large lepton asymmetry
($\xi_\nu\gtrsim10^{-2}$) would require a lepton-number source acting \emph{after}
sphaleron freeze-out---which inheritance-then-sphaleron-processing does not
supply---and would disfavour the framework's account of the asymmetry. A null
result ($\xi_\nu$ consistent with $\sim\eta$) is a weak confirmation. Improved
BBN$+$CMB analyses tighten $\xi_\nu$ now; eventual direct detection of the cosmic
neutrino background (PTOLEMY-class experiments) would additionally probe the
handedness content---a pure-continuation lineage carries our matter neutrinos,
with no antineutrino-like excess from a $CPT$ mirror. Reproducible as
\texttt{sims/src/cartasis\_sims/lepton\_asymmetry.py}.

\section{Primordial gravitational waves}

Inflation predicts a tensor-to-scalar ratio $r$ in a fairly definite range; a
torsion bounce, with its different early-expansion history
(Chapter~\ref{chap:bounce}), predicts a different and typically \emph{smaller}
$r$. This is a clean fork. LiteBIRD will constrain $r$ to the $\sim\!10^{-3}$
level: a detection at the higher values favours inflation, while a value pushed
low or absent favours the bounce. Turning ``typically smaller'' into a definite
predicted band is a computation the framework owes (Chapter~\ref{chap:sims}), but
the \emph{direction} of the discriminator is already set.

\section{The dark-matter null}

Because dark matter here is geometric---the parent's projected mass, not a
particle (Chapter~\ref{chap:home})---the framework predicts that direct-detection
experiments will keep finding \emph{nothing}. WIMPs, axions, sterile neutrinos:
the prediction is a continued null across LZ, XENONnT~\citep{aprile2023}, and
their successors as
they probe ever-lower cross-sections. This is a weak test taken one experiment at a
time---a null is always consistent with a particle just out of reach---but its
weight accumulates, and it is a prediction the framework makes by construction
while particle dark matter must keep explaining the silence away.

\section{Large-scale structure}

If dark matter traces the parent's projected mass rather than a clustering
particle, the relationship between the baryon acoustic oscillations and the
dark-matter distribution should follow patterns specific to the framework, and the
largest scales---where the inherited lumpiness of Chapter~\ref{chap:home} lives---
should show departures from $\Lambda$CDM. DESI, Euclid, and the Roman Space
Telescope will map structure with the precision to see them, and the prediction is
pointed: excess large-scale power, correlated with the dark background rather than
the luminous one, on scales above our own homogeneity scale.

\medskip
\noindent Of these, the cleanest single forks are the tensor ratio $r$ (a bounce
runs low where inflation can run high) and the lepton asymmetry (the framework
forbids the large value the data still allow). The dark-matter null and the
large-scale-structure patterns accumulate weight more slowly but in one
direction. Each is logged, with its current status and the calculation it still
needs, in the scoreboard of Chapter~\ref{chap:scoreboard}.
